\chapter{Zanurzenia i otoczenia podrozmaitości}
\label{ch:wlozenia-otoczenia}
\index{zanurzenie@Zanurzenie}

W Rozdziale~\ref{ch:rozmaitosci} poznaliśmy zanurzenia
(ang.\ \emph{embeddings}) i mapy plasterkowe.
Teraz pytamy o całe \emph{otoczenie} podrozmaitości:
czy można je opisać za pomocą wektorów wychodzących
z podrozmaitości? Odpowiedź daje wiązka normalna.
Jej lokalny obraz przypomina wąski pas wokół krzywej;
globalnie pas może być skręcony, jak wstęga Möbiusa.
To rozróżnienie będzie istotne przy dołączaniu uchwytów
i operacjach chirurgicznych.

Zależności między wynikami tego rozdziału są następujące.
Twierdzenie Whitneya umieszcza abstrakcyjną rozmaitość
w przestrzeni euklidesowej. Wiązka normalna opisuje
kierunki poprzeczne do tak umieszczonej rozmaitości,
a twierdzenie o otoczeniu tubularnym zamienia krótkie
wektory normalne w rzeczywiste punkty otoczenia.
Kołnierz robi podobną rzecz dla brzegu. Funkcje
wyczerpujące pozwalają stosować te konstrukcje także
na rozmaitościach niezwartych.

\section{Zanurzenia właściwe i twierdzenie Whitneya}
\label{sec:whitney-wlozenia}

Przez \emph{zanurzenie}, zdefiniowane już
w Rozdziale~\ref{ch:rozmaitosci}, rozumiemy immersję,
która jest homeomorfizmem na obraz. Nie należy go
mylić z samą różnowartościową immersją;
Przykład~\ref{prz:osemka} pokazuje różnicę.
W tym rozdziale $M$ i $N$ są gładkimi rozmaitościami
Hausdorffa z przeliczalną bazą i bez brzegu, chyba że
zaznaczymy inaczej (jak w części o kołnierzach).

\begin{definicja}[Odwzorowanie właściwe]
\index{odwzorowanie wlasciwe@Odwzorowanie właściwe}
\label{def:odwzorowanie-wlasciwe}
Odwzorowanie ciągłe $F:M\to N$ jest \emph{właściwe},
jeżeli przeciwobraz każdego zbioru zwartego w $N$
jest zwarty w $M$. Zanurzenie właściwe to zanurzenie
będące odwzorowaniem właściwym.
\end{definicja}

\begin{stwierdzenie}[Kiedy immersja jest zanurzeniem]
\label{prop:proper-immersja}
Różnowartościowa immersja właściwa $F:M\to N$
jest zanurzeniem, a jej obraz jest domknięty w $N$.
Jeśli $M$ jest zwarta, wystarczy, aby immersja
była różnowartościowa.
\end{stwierdzenie}
\begin{proof}
Pokażemy najpierw, że $F$ jest odwzorowaniem
domkniętym. Niech $A\subset M$ będzie domknięty
i niech $y_j=F(x_j)\in F(A)$ zbiega do $y\in N$.
Zbiór $C=\{y\}\cup\{y_j:j\geq1\}$ jest zwarty
(otwarte otoczenie $y$ zawiera prawie wszystkie
wyrazy ciągu). Z właściwości $F^{-1}(C)$ jest
zwarty. Ciąg $x_j$ ma więc podciąg zbieżny
do $x\in F^{-1}(C)$. Ponieważ $A$ jest domknięty,
$x\in A$, a ciągłość daje $F(x)=y$. Rozmaitości
są metryzowalne, toteż wykazane domknięcie
sekwencyjne jest domknięciem topologicznym.

Różnowartościowe odwzorowanie ciągłe i domknięte
jest homeomorfizmem na obraz: dla każdego
domkniętego $A\subset M$ jego obraz jest
domknięty również w $F(M)$, co jest właśnie
ciągłością $F^{-1}:F(M)\to M$. W połączeniu
z założeniem o immersji daje to zanurzenie.
Biorąc $A=M$, otrzymujemy domknięcie $F(M)$
w $N$. Gdy $M$ jest zwarta, ciągłe $F$
jest właściwe, bo przeciwobrazy zwartych
zbiorów są domkniętymi podzbiorami $M$.
\end{proof}

\begin{uwaga}
Zanurzenie nie musi być właściwe:
$(0,1)\hookrightarrow\R$ jest zanurzeniem,
ale jego obraz nie jest domknięty.
W twierdzeniu tubularnym dla niezwartych
podrozmaitości domknięcie obrazu daje wygodną
kontrolę globalną.
\end{uwaga}

\begin{twierdzenie}[Whitney: zanurzenie do przestrzeni euklidesowej]
\label{thm:whitney-slabe}
Każda zwarta gładka rozmaitość $M^n$ bez brzegu
ma gładkie zanurzenie do $\R^{2n+1}$.
\end{twierdzenie}
\begin{proof}
\emph{Krok 1: zanurzenie do pewnej $\R^N$.}
Wybierzmy skończone pokrycie mapami $(U_i,\varphi_i)$,
$i=1,\ldots,r$, o ograniczonych obrazach. Takie pokrycie
otrzymujemy, wybierając przy każdym punkcie mniejszą
kulę współrzędnych i korzystając ze zwartości. Niech $\chi_i$ będą gładkimi
funkcjami podziału jedności, z
$\supp\chi_i\subset U_i$. Przez zero poza $U_i$
przedłużamy funkcje
$\widetilde\varphi_i=\chi_i\varphi_i$.
Przedłużenie jest gładkie: przy punkcie poza
$\supp\chi_i$ obie funkcje są równe zeru
w pewnym otoczeniu. Definiujemy
\[
 F(x)=\bigl(\chi_1(x),\widetilde\varphi_1(x);
       \ldots;\chi_r(x),\widetilde\varphi_r(x)\bigr)
       \in\R^{r(n+1)}.
\]
Jeżeli $F(x)=F(y)$, wybierzmy $i$ z
$\chi_i(x)>0$; istnieje taki $i$, bo
$\sum_i\chi_i(x)=1$. Z pierwszej składowej
bloku otrzymujemy $\chi_i(y)=\chi_i(x)>0$,
więc $x,y\in U_i$. Z następnych składowych
otrzymujemy
$\chi_i(x)\varphi_i(x)=\chi_i(y)\varphi_i(y)$,
a zatem $\varphi_i(x)=\varphi_i(y)$ i $x=y$.
Jeżeli $\dd F_x(v)=0$, składowa $\chi_i$
daje $\dd\chi_i{}_x(v)=0$. Różniczkując
$\chi_i\varphi_i$, dostajemy
\[
 0=\dd(\chi_i\varphi_i)_x(v)
 =\dd\chi_i{}_x(v)\varphi_i(x)
       +\chi_i(x)\dd\varphi_i{}_x(v)
 =\chi_i(x)\dd\varphi_i{}_x(v).
\]
Ponieważ $\chi_i(x)>0$ i $\dd\varphi_i{}_x$
jest izomorfizmem, $v=0$. Tak więc $F$
jest różnowartościową immersją, a ze
zwartości $M$ i Stwierdzenia~\ref{prop:proper-immersja}
wynika, że jest zanurzeniem.

\emph{Krok 2: zmniejszanie wymiaru.}
Możemy dopisać zerowe współrzędne, aby
założyć $N>2n+1$. Dla niezerowego
$v\in\R^N$ niech $P_v$ będzie rzutem
ortogonalnym na hiperpłaszczyznę
$v^\perp\cong\R^{N-1}$. Rzut $P_v\circ F$
przestaje być różnowartościowy dokładnie
wtedy, gdy prosta $\R v$ jest kierunkiem
siecznej $F(x)-F(y)$ dla pewnych $x\ne y$.
Zbiór tych kierunków jest obrazem
gładkiego odwzorowania
\[
 s:(M\times M)\setminus\Delta\longrightarrow
       \R P^{N-1},\qquad
 s(x,y)=[F(x)-F(y)].
\]
Dziedzina ma wymiar $2n$, a $\R P^{N-1}$
ma wymiar $N-1>2n$. Żadna różniczka $s$
nie może więc być surjektywna.
Z twierdzenia Sarda~\ref{thm:sard} obraz
$s$ ma miarę zero w mapach przestrzeni
projektowej.

Trzeba również zachować immersję. Jej
różniczka straciłaby injektywność,
gdyby $\R v$ była obrazem pewnego
niezerowego wektora stycznego.
Takie kierunki są obrazem odwzorowania
\[
 t:\mathbb P(TM)\longrightarrow\R P^{N-1},
       \qquad t(x,[w])=[\dd F_x(w)].
\]
Przestrzeń prostych stycznych
$\mathbb P(TM)$ jest lokalnie iloczynem
mapy na $M$ z $\R P^{n-1}$, a zatem
ma wymiar $n+(n-1)=2n-1$.
Również obraz $t$ ma miarę zero.
Jeżeli $n=0$, zbiór prostych stycznych
jest pusty i ten warunek znika.

Dwa zbiory miary zero nie pokrywają
$\R P^{N-1}$. Wybierzmy prostą $[v]$
poza ich sumą. Wtedy $P_v\circ F$ jest
różnowartościową immersją, a więc
zanurzeniem na mocy zwartości $M$.
Powtarzamy redukcję, dopóki wymiar
przestrzeni docelowej nie wyniesie
$2n+1$. W każdym kroku przesłanka
$N>2n+1$ jest zachowana do chwili
ostatniego rzutowania.
\end{proof}

\begin{twierdzenie}[Mocne twierdzenie Whitneya]
\label{thm:whitney-mocne}
Każda zwarta gładka rozmaitość $M^n$ bez brzegu
o $n\geq1$ ma gładkie zanurzenie w $\R^{2n}$.
\end{twierdzenie}

W dowodzie nie można po prostu powtórzyć
poprzedniego rzutowania: gdy obraz ma już
wymiar $2n+1$, zbiór kierunków siecznych
może wypełnić dodatnią część przestrzeni
kierunków. Rzut do $\R^{2n}$ będzie jednak
zazwyczaj \emph{immersją z izolowanymi
punktami podwójnymi}. Trzeba następnie
usunąć te punkty. Wyjaśnimy osobno
konstrukcję geometryczną, która to umożliwia.

\begin{stwierdzenie}[Dysk Whitneya i usuwanie pary]
\label{lem:dysk-whitneya-15}
Niech $n\geq3$, niech $M$ będzie zwarta i niech
$f:M^n\to\R^{2n}$ będzie immersją z poprzecznymi punktami
podwójnymi, bez punktów potrójnych.
Przypuśćmy, że dwie takie pary mają
preobrazy $x_1,y_1$ oraz $x_2,y_2$,
a wybór łuków $x_1x_2$ i $y_1y_2$
jest \emph{zgodny} w tym sensie:
dwa lokalne znaki przecięcia przy końcach
są przeciwne po przeniesieniu orientacji
wzdłuż łuków. Wtedy istnieje modyfikacja
$f$, wsparta w otoczeniu obu punktów
i łączącego je dysku, która usuwa
te dwa punkty podwójne i nie tworzy nowych.
W przypadku zorientowanego $M$ i parzystego
$n$ zgodność oznacza po prostu przeciwne
znaki przecięć.
\end{stwierdzenie}
\begin{proof}
Punkty podwójne mają skończenie wiele preobrazów.
Istotnie, transwersalność izoluje pary $(x,y)$ poza
przekątną $M\times M$, a lokalna różnowartościowość
immersji wyklucza ich gromadzenie się przy przekątnej.
Zwartość daje skończoność.
Przy $n\geq3$ możemy wybrać łuki w $M$,
które nie spotykają innych preobrazów
punktów podwójnych i są rozłączne:
dwie krzywe w $n$ wymiarach po małej perturbacji
się mijają, bo $1+1<n$. Niewielka
zmiana łuków zapewnia, że ich obrazy
przez $f$ także są rozłączne poza
obrazami końców. W obrazie tworzą
prostą pętlę z dwoma narożami.
Ponieważ $\R^{2n}$ jest po prostu spójna,
pętla ogranicza odwzorowanie dysku. Najpierw przy
jej brzegu ustalamy wąski pas, którego wnętrze
omija $f(M)$; przy narożach robimy to w modelu
dwóch poprzecznych płaszczyzn. Wypełnienie można
wygładzić względnie tego pasa (współrzędne są
rzeczywiste, więc wystarczy lokalne wygładzanie
i podział jedności). Dysk ma dwa naroża na brzegu;
jego powierzchnia nośna jest gładka także przy nich.

Uzasadnijmy dokładniej użyte tu położenie ogólne.
W małej mapie wnętrza o współrzędnych $(u_1,u_2)$
dodajemy do wypełnienia perturbacje
$\beta(u)(a_0+a_1u_1+a_2u_2)$, gdzie $a_i\in\R^{2n}$,
a $\beta=1$ na mniejszej mapie i znika poza większą.
Parametry pozwalają niezależnie zmieniać wartość
i obie kolumny różniczki. Warstwa macierzy
$2n\times2$ rzędu $r<2$ ma kowymiar
$(2n-r)(2-r)\geq2n-1>2$. Wzór na kowymiar wynika
z eliminacji blokowej przy odwracalnym minorze
rzędu $r$: pozostały blok jest wyznaczony przez
pozostałe trzy bloki. Parametryczna transwersalność
do tych warstw zapewnia więc, że mała perturbacja
jest immersją. Skończone pokrycie zwartej części
dysku poza ustalonym pasem daje skończenie wiele
parametrów; przy brzegu immersja jest już ustalona.

Dla otrzymanej immersji pary dostatecznie bliskich
punktów nie dają samoprzecięć, także po małej zmianie
w $C^1$. Można to sprawdzić na skończonym pokryciu
mapami, w których odpowiedni dwuwymiarowy rzut ma
różniczkę bliską ustalonemu izomorfizmowi.
Pozostałe możliwe pary leżą w zwartym zbiorze
oddalonym od przekątnej. Na rozłącznych mapach
dwóch argumentów funkcje odcinające pozwalają
zmieniać wartości niezależnie. Transwersalność
różnicy wartości do $0\in\R^{2n}$ usuwa te pary,
gdyż $2+2-2n<0$. Tak samo dla pary składającej się
z punktu wnętrza dysku i punktu $M$ zmieniamy
wartość po stronie dysku; wymiar potencjalnego
przecięcia to $2+n-2n<0$. W ustalonym pasie nie
ma niepożądanych przecięć, a jego węższy podpas
możemy pozostawić nietknięty. Zwartość pozwala
wybrać skończone rodziny perturbacji na pozostałym
obszarze i zachować uzyskaną wcześniej immersję.
Otrzymujemy zanurzony dysk $D$, którego wnętrze
nie przecina $f(M)$.

Pozostaje sprawdzić, czy można wykonać
ruch jednego płata \emph{przez} ten dysk,
nie skręcając go po drodze. Wiązka
normalna dysku ma rząd $2n-2$ i jest
trywialna, ponieważ dysk jest ściągalny.
Oznaczmy $k=n-1\geq2$. Na pierwszym łuku normalne
do łuku wewnątrz pierwszego płata dają $k$-płaszczyznę
$A$ w $\nu D$, a na drugim analogicznie dostajemy $B$.
Szukamy rozkładu $\nu D=\xi_1\oplus\xi_2$, w którym
$\xi_1=A$ na pierwszym łuku, zaś $\xi_2=B$ na drugim.
Po wybraniu pomocniczej metryki i standardowego modelu
przy narożach odpowiada to pętli $L$ w
$\operatorname{Gr}_k(\R^{2k})$: na pierwszym łuku
przepisujemy $L=A$, a na drugim $L=B^\perp$.
Przy narożach płaszczyzny $A$ i $B$ są dopełniające;
metrykę można tam dobrać tak, by były ortogonalne.
Ważne jest użycie dopełnienia na drugim łuku.

Pętla $L$ przedłuża się na dysk dokładnie wtedy,
gdy zachowuje orientację swojej $k$-płaszczyzny.
Oto uzasadnienie. Przestrzeń zorientowanych płaszczyzn
\[
 \operatorname{Gr}^+_k(\R^{2k})
 =\operatorname{SO}(2k)/(\operatorname{SO}(k)\times
                         \operatorname{SO}(k))
\]
jest po prostu spójna. Wyjaśnijmy używaną własność
wiązania: każdą pętlę w ilorazie można podnieść do
drogi w grupie, a jej końce połączyć we włóknie,
ponieważ $\operatorname{SO}(k)\times\operatorname{SO}(k)$
jest spójna. Dostajemy pętlę w grupie o zadanym rzucie
(z dokładnością do dodania stałego odcinka).
Jeżeli każdą taką pętlę można zdeformować do pętli
w samym włóknie, jej rzut jest ściągalny.
Wystarczy zatem sprawdzić surjektywność
$\pi_1(\operatorname{SO}(k)\times\operatorname{SO}(k))
\to\pi_1\operatorname{SO}(2k)$: obrót o $2\pi$
w jednej ustalonej dwupłaszczyźnie pierwszego czynnika
daje generator grupy docelowej $\Z/2\Z$.
Przypomnijmy uzasadnienie tego ostatniego faktu:
$\operatorname{SO}(3)$ ma dwukrotne nakrycie przez
jednostkowe kwaterniony $S^3$, a obrót o $2\pi$
podnosi się do drogi od $1$ do $-1$.
Wiązanie $\operatorname{SO}(r-1)\to\operatorname{SO}(r)
\to S^{r-1}$ wysyła macierz na jej ostatnią kolumnę.
Dla $r\geq4$ sfera ma zerowe $\pi_1$ i $\pi_2$:
przez aproksymację komórkową każde odwzorowanie z okręgu
lub dwusfery jest homotopijne z odwzorowaniem stałym.
Podnosząc ściągnięcie rzutu pętli, deformujemy ją
do włókna. Jeśli pętla we włóknie ogranicza dysk
w całej grupie, rzut dysku z zapadniętym brzegiem
jest odwzorowaniem $S^2\to S^{r-1}$. Jego ściągnięcie, podniesione
względnie brzegu, deformuje dysk do włókna.
Inkluzja włókna daje więc izomorfizm grup podstawowych
i przenosi wynik z $r=3$ na wszystkie $r\geq4$.
Użyte podnoszenia wynikają z lokalnych trywializacji:
dzielimy przedział lub kwadrat parametrów na małe
kawałki i w każdym podnosimy przez iloczyn.
Trywializacje obu wiązań otrzymujemy, lokalnie
uzupełniając ortonormalne bazy metodą Grama--Schmidta.
Dla $k=2$ czynnik
$\operatorname{SO}(2)$ ma grupę podstawową $\Z$,
ale ten sam obrót nadal daje generator obrazu.
Dwukrotne nakrycie przez zorientowany grassmannian
pokazuje więc, że $\pi_1\operatorname{Gr}_k(\R^{2k})
=\Z/2\Z$, a jego jedyną przeszkodą jest odwrócenie
orientacji. Reprezentantem klasy niezerowej jest
\[
 L_\theta=\operatorname{span}
 (\cos\theta\,e_1+\sin\theta\,e_{k+1},e_2,\ldots,e_k),
 \qquad 0\leq\theta\leq\pi.
\]

Porównajmy teraz znaki. Orientacje obu płatów
przenosimy wzdłuż wybranych łuków. W bazie płata
oddzielamy kierunek łuku od pozostałych $k$ kierunków.
Przy dwóch końcach para kierunków łuków wyznacza
przeciwne orientacje płaszczyzny dysku: łuki
ograniczają dwukąt. Aby całe bazy otoczenia miały
przeciwne znaki przecięcia, pozostałe kierunki
muszą więc skleić się bez odwrócenia orientacji.
Jest to właśnie trywialność przeszkody dla $L$.
Zgodność znaków daje zatem przedłużenie $L$ na $D$,
a jego dopełnienie daje $\xi_2$. Obie wiązki nad
dyskiem są trywialne, więc teraz wybieramy ich ramki.
Przedłużaliśmy rozkład na podwiązki; nie twierdzimy,
że każda z góry zadana pełna ramka na brzegu
przedłuża się na dysk.

W tych ramkach otoczenie dysku ma współrzędne
$D\times\R^k\times\R^k$ o dostatecznie małym promieniu.
Użycie twierdzenia tubularnego jest tu niezależne
od twierdzenia Whitneya: jego dowód poniżej korzysta
tylko z metryki, funkcji wyczerpującej i twierdzenia
odwrotnego. Przy brzegu i narożach dysku najpierw
nieco przedłużamy jego powierzchnię nośną, a potem
ograniczamy otoczenie do zwartego $D$.
Prostując płaty w wąskich pasach przy łukach,
otrzymujemy pierwszy łuk pogrubiony w pierwszym
$\R^k$ i drugi w drugim $\R^k$.

Sam ruch w dwuwymiarowym przekroju można opisać
jawnie. Po wyborze współrzędnych łuki mają postać
$y=x(1-x)$ i $y=-x(1-x)$ dla $0\leq x\leq1$,
z krótkimi przedłużeniami poza końce. Pierwszy
zastępujemy rodziną wykresów
$y=x(1-x)-s h(x)$, $0\leq s\leq1$, gdzie $h\geq0$
ma nośnik w przedłużonym przedziale i
$h(x)>2x(1-x)$ na $[0,1]$.
Na końcu pierwszy wykres leży ściśle poniżej
drugiego; poza przedziałem już wcześniej był poniżej.
Każdy wykres pozostaje regularny. Pogrubiając go
w kierunkach $\xi_1$, przenosimy ten ruch na płat;
kierunki $\xi_1$ i $\xi_2$ są dopełniające, więc
usunięcie obu przecięć wykresów usuwa też przecięcia
płatów. Ruch wygaszamy przy zewnętrznym brzegu
pogrubienia. Wnętrze dysku omija $f(M)$, a zwartość
zapewnia, że poza ustalonymi pasami możemy wybrać
otoczenie nieprzecinające pozostałej części obrazu.
Przy końcach używamy modelu dwóch poprzecznych
płaszczyzn. Powstała immersja nie ma nowych
punktów podwójnych.
\end{proof}

\begin{figure}[htbp]
\centering
\begin{tikzpicture}[>=Latex,scale=.95]
  \fill[manifoldgold!13]
    plot[domain=0:3.2,samples=65] (\x,{.35*\x*(3.2-\x)})
    -- plot[domain=3.2:0,samples=65] (\x,{-.35*\x*(3.2-\x)}) -- cycle;
  \draw[manifoldblue!85!black,very thick]
    plot[domain=-.35:3.55,samples=81] (\x,{.35*\x*(3.2-\x)});
  \draw[manifoldred!85!black,very thick]
    plot[domain=-.35:3.55,samples=81] (\x,{-.35*\x*(3.2-\x)});
  \fill[manifoldblue!75!black] (0,0) circle (2.2pt)
    node[left=5pt] {$z_1$};
  \fill[manifoldblue!75!black] (3.2,0) circle (2.2pt)
    node[right=5pt] {$z_2$};
  \node[manifoldgold!75!black] at (1.6,0) {$D$};
  \draw[->,manifoldgold!80!black,thick] (1.6,1.55)--(1.6,.95)
    node[pos=0,above] {przesunięcie płata};
\end{tikzpicture}
\caption{Schemat dysku Whitneya $D$: dwa łuki jego brzegu
leżą na różnych płatach immersji i łączą dwa punkty
podwójne $z_1,z_2$. Ruch płata po dysku usuwa oba
przecięcia, gdy obramowanie brzegowe da się
przedłużyć na $D$. Rysunek pokazuje tylko
dwuwymiarowy przekrój konstrukcji.}
\label{fig:dysk-whitneya-15}
\end{figure}

\begin{proof}[Dowód Twierdzenia~\ref{thm:whitney-mocne}]
Najpierw zakładamy $n\geq3$ i $M$ spójną.
Z Twierdzenia~\ref{thm:whitney-slabe}
mamy zanurzenie $F:M\hookrightarrow\R^{2n+1}$.
Rzutowanie w kierunku $v$ daje
$f=P_v\circ F:M\to\R^{2n}$. Kierunki,
dla których $f$ nie jest immersją,
tworzą obraz $\mathbb P(TM)$ o wymiarze
$2n-1$ w przestrzeni kierunków
$\R P^{2n}$: podobnie jak w poprzednim
dowodzie można więc wybrać $v$ poza
tym obrazem.

Rozważmy odwzorowanie kierunków siecznych
\[
 s:(M\times M)\setminus\Delta_M
   \longrightarrow\R P^{2n},\qquad
 s(x,y)=[F(x)-F(y)].
\]
Twierdzenie Sarda daje regularny kierunek
$[v]$, który jednocześnie omija
kierunki styczne. Wtedy
$f(x)=f(y)$, $x\ne y$, oznacza
$s(x,y)=[v]$. Regularność $s$
oznacza, że odwzorowanie
\[
 T_xM\oplus T_yM\longrightarrow\R^{2n},
 \qquad (u,w)\longmapsto
   \dd f_x(u)-\dd f_y(w)
\]
jest surjektywne, czyli obrazy dwóch
$n$-płaszczyzn stycznych są poprzeczne.
Rozwiązania są izolowane. Nie mogą
gromadzić się przy przekątnej:
granice kierunków bardzo krótkich
siecznych byłyby kierunkami stycznymi,
a wybraliśmy $[v]$ poza nimi. Dokładniej, w jednej
mapie dzielimy różnicę $F(x_j)-F(y_j)$ przez
odległość współrzędnych argumentów. Po wybraniu
podciągu jednostkowych kierunków argumentów wzór
całkowy na przyrost daje w granicy $\dd F_x(w)$
dla pewnego $|w|=1$. Immersja zapewnia, że wektor
ten jest niezerowy, co daje sprzeczność.
Zwartość $M$ daje zatem skończenie
wiele punktów podwójnych.
Małą perturbacją $f$ usuwamy też
ewentualne punkty potrójne.
W istocie dla trzech różnych argumentów
warunki $f(x)=f(y)$ i $f(x)=f(z)$
mają razem $4n$ równań, podczas gdy
$(x,y,z)$ ma wymiar $3n$; perturbacje
funkcjami o rozłącznych nośnikach
przy $x,y,z$ pozwalają zmieniać
obie różnice niezależnie.
Transwersalność z parametrem
(Twierdzenie~\ref{thm:transwersalnosc-param})
daje więc brak punktów potrójnych. Aby zastosować
ją do jednej skończonej rodziny, ustalamy
$\delta>0$ tak małe, że każda dostatecznie bliska
w $C^1$ perturbacja jest różnowartościowa na parach
o odległości mniejszej niż $\delta$. Istnienie
$\delta$ wynika z lokalnych rzutów immersji na
współrzędne i skończonego pokrycia zwartej $M$.
Wszystkie możliwe trójki leżą więc w zwartym zbiorze
o wzajemnych odległościach argumentów co najmniej
$\delta$. Pokrywamy go skończenie wieloma produktami
trzech rozłącznych obszarów map. Stałe przesunięcia
pomnożone przez funkcje odcinające, równe $1$ na
mniejszych obszarach, dają rodzinę o surjektywnej
pochodnej parametrycznej obu różnic. Transwersalność
na otoczeniu tego zwartego zbioru usuwa wszystkie
trójki jednocześnie. Dostateczna małość perturbacji
zachowuje immersję oraz poprzeczność par: potencjalne
pary też leżą w zwartym zbiorze poza przekątną,
gdzie stosujemy otwartość transwersalności.

Pokażemy, skąd wziąć partnera dla
punktu podwójnego o niezgodnym znaku.
Istnieje jawna immersja sfery
$S^n=\{(u,x)\in\R\times\R^n:
u^2+|x|^2=1\}$ w $\R^{2n}$:
\[
 W(u,x)=(ux,x).
\]
Jedyny jej punkt podwójny to $0$,
o preobrazach $(1,0)$ i $(-1,0)$.
Przy tych punktach styczne obrazy
mają postać $\{(h,h)\}$ oraz
$\{(-h,h)\}$, więc są poprzeczne.
Gdy $x\ne0$, druga składowa
wyznacza $x$, pierwsza wyznacza $u$;
zerowanie różniczki wymusza najpierw $\delta x=0$,
a następnie $x\delta u=0$, więc $\delta u=0$.
Przy $x=0$ warunek styczności do sfery daje
$\delta u=0$, a druga składowa różniczki wykrywa
każde $\delta x\ne0$. Jest to zatem immersja wszędzie. Umieszczamy małą
kopię tej immersji w wolnym
obszarze i łączymy jej regularny
fragment z regularnym fragmentem
$f(M)$ wąską rurką. Ponieważ
$1+n<2n$, rdzeń rurki można
poprowadzić poza pozostałym
obrazem; otrzymujemy immersję
$M\#S^n\cong M$ z jednym
dodatkowym punktem podwójnym. Używamy standardowego
sklejenia brzegów usuniętych dysków: dołączona
sfera z usuniętym dyskiem jest zwykłym dyskiem,
więc nie zmieniamy gładkiego typu $M$.
Odbicie jednej współrzędnej
otoczenia odwraca jego znak
w przypadku parzystego
$n$ i zorientowanego $M$.
Przy nieparzystym $n$ zamiana
ról dwóch lokalnych płatów
odwraca znak względny,
a gdy $M$ jest niezorientowana,
dobór łuku lub jego obieg wokół
pętli odwracającej orientację
zmienia względne obramowanie.
Możemy więc dodać punkt
zgodny z danym w sensie
Stwierdzenia~\ref{lem:dysk-whitneya-15}.
Po usunięciu pary liczbę
punktów podwójnych zmniejszamy
o jeden. Ponieważ początkowo
było ich skończenie wiele,
powtarzanie daje różnowartościową
immersję. Zwartość $M$ czyni
ją zanurzeniem.

Dla $n=1$ każda zwarta spójna
rozmaitość jest okręgiem;
rozłączne składowe umieszczamy
jako rozłączne okręgi w $\R^2$.
Dla $n=2$ poprzedni ruch nie
działa: $2+2-4=0$, więc dysk
Whitneya może przecinać
zarówno siebie, jak i powierzchnię.
Korzystamy z klasyfikacji
zwartych powierzchni:
zorientowane powstają jako
sumy spójne torusów i sfer,
a niezorientowane jako sumy
spójne płaszczyzn projektowych.
Zorientowane realizujemy w $\R^3$.
Dla $\R P^2$ mamy jawne zanurzenie
do $\R^4$. Na $S^2$ połóżmy
$w=x+iy$ i
\[
 V([x:y:z])
   =\bigl(\operatorname{Re}(w^2),
           \operatorname{Im}(w^2),
           \operatorname{Re}(zw),
           \operatorname{Im}(zw)\bigr).
\]
Wartość nie zmienia się przy
$(x,y,z)\mapsto(-x,-y,-z)$,
więc wzór schodzi na $\R P^2$. Zapis $[x:y:z]$
oznacza tu wybór reprezentanta jednostkowego;
dla dowolnego niezerowego reprezentanta wszystkie
cztery składowe należy podzielić przez $x^2+y^2+z^2$.
Jeśli $w\ne0$, równość pierwszych
dwóch składowych daje $w'=\pm w$;
pozostałe dają $z'=\pm z$
z tym samym znakiem. Jeśli $w=0$,
obie możliwości są biegunami
identyfikowanymi w $\R P^2$.
Jest więc różnowartościowy.
Ponadto zerowanie różniczki
pierwszych dwóch składowych
przy $w\ne0$ daje $\delta w=0$,
a ostatnich dwóch $\delta z=0$.
Przy $w=0$ mamy $z=\pm1$;
zerowanie ostatnich dwóch
również daje $\delta w=0$,
a warunek styczności do $S^2$
daje $\delta z=0$. Jest immersją,
więc ze zwartości jest zanurzeniem.
Umieszczając rozłączne powierzchnie
w osobnych kulach $\R^4$ i łącząc
je cienkimi rurkami wzdłuż
rozłącznych łuków w dopełnieniu
($1+2<4$), realizujemy ich
sumy spójne. Tak otrzymujemy
wszystkie zwarte powierzchnie.
Składowych zwartej rozmaitości jest skończenie
wiele, ponieważ są otwarte i tworzą pokrycie.
Wreszcie składowe rozmaitości $M$ umieszczamy w rozłącznych
kulach $\R^{2n}$.
\end{proof}

\begin{uwaga}[Zakres twierdzenia]
Pełne mocne twierdzenie Whitneya
obejmuje też odpowiednie rozmaitości
niezwarte. Powyższy dowód dotyczy
zwartych rozmaitości bez brzegu,
zgodnie z zakresem wcześniejszego
Twierdzenia~\ref{thm:whitney-slabe}.
W wymiarze dwa użyliśmy
klasyfikacji powierzchni; nie
podstawiliśmy błędnie argumentu
z dyskiem Whitneya w wymiarze cztery.
Twierdzenie dotyczy gładkiej struktury,
nie zachowania zadanej metryki:
zanurzenie $M\to\R^{2n}$ nie musi
być izometryczne.
\end{uwaga}

\begin{przyklad}
Okrąg $S^1$ jest zanurzony w $\R^2$ przez
$\theta\mapsto(\cos\theta,\sin\theta)$
(ściśle: jest to odwzorowanie z $\R/2\pi\Z$).
Wersja $2n+1$ gwarantuje jedynie
zanurzenie w $\R^3$: jest twierdzeniem
uniwersalnym, nie wzorem na najmniejszy
wymiar dla konkretnej rozmaitości.
Podobnie $S^2\hookrightarrow\R^3$ ma
niższy wymiar niż ogólna gwarancja $\R^5$.
\end{przyklad}

\section{Wiązka normalna: kierunki poprzeczne}
\label{sec:wiazka-normalna}

Niech $i:S^k\hookrightarrow M^n$ będzie
zanurzeniem. W każdym $x\in S$ różniczka
$\dd i_x$ identyfikuje $T_xS$ z
podprzestrzenią $T_{i(x)}M$.
Pisząc $S\subset M$, pomijamy $i$ w zapisie.

\begin{definicja}[Wiązka normalna]
\index{wiazka normalna@Wiązka normalna}
\label{def:wiazka-normalna-15}
\emph{Wiązką normalną} zanurzenia $S\subset M$
nazywamy wiązkę ilorazową
\[
 \nu(S,M):=TM|_S/TS,\qquad
 \nu_x(S,M)=T_xM/T_xS.
\]
Jej rząd to $n-k$, czyli kowymiar $S$.
Krótki ciąg wiązek
\begin{equation}\label{eq:normalny-ciag}
 0\longrightarrow TS
 \stackrel{\dd i}{\longrightarrow}TM|_S
 \longrightarrow\nu(S,M)\longrightarrow0
\end{equation}
jest dokładny we włóknie nad każdym $x$.
\end{definicja}

Klasa $[v]\in T_xM/T_xS$ pamięta tę część
przemieszczenia, której nie da się osiągnąć
ruchem stycznym po $S$. Wybór metryki
riemannowskiej $g$ wskazuje jednego
przedstawiciela każdej klasy: wektor
prostopadły do $T_xS$.

\begin{stwierdzenie}[Gładkość i rozszczepienie]
\label{prop:normalna-rozszczepienie}
Wiązka $\nu(S,M)$ jest gładką wiązką
wektorową rzędu $n-k$. Każda metryka
riemannowska $g$ na $M$ daje izomorfizm
wiązki $\nu(S,M)\cong(TS)^\perp\subset TM|_S$
oraz rozkład $TM|_S=TS\oplus(TS)^\perp$.
\end{stwierdzenie}
\begin{proof}
W mapie plasterkowej $S$ jest zbiorem
$\{(u,0):u\in\R^k\}$, a wektor w $T_{(u,0)}M$
ma współrzędne $(a,b)\in\R^k\oplus\R^{n-k}$.
Klasy modulo $TS$ są dokładnie parami
$(u,b)$. Zmiana mapy daje gładką zmianę
baz we włóknach, dlatego lokalne opisy
sklejają się w wiązkę gładką.

Niech $P^\perp_x:T_xM\to(T_xS)^\perp$
będzie rzutem ortogonalnym.
Jądrem $P^\perp_x$ jest $T_xS$, więc
odwzorowanie $[v]\mapsto P^\perp_xv$ jest
dobrze określonym izomorfizmem.
W lokalnej ramie $TS$ współczynniki
rzutu otrzymuje się przez odwrócenie
macierzy Grama
$(g(e_a,e_b))_{a,b=1}^k$;
macierz ta jest odwracalna i gładko
zależy od $x$. Rzut, a więc również
izomorfizm wiązek, jest gładki.
\end{proof}

\begin{przyklad}[Dwa przeciwstawne pasy normalne]
\label{prz:normalna-mobius}
Dla $S^1\subset\R^2$ wektor
$n(\theta)=(\cos\theta,\sin\theta)$
jest globalnym niezerowym polem normalnym.
Wiązka normalna jest iloczynem
$S^1\times\R$.
Rozważmy teraz środkowy okrąg wstęgi
Möbiusa
\[
 M_{\mathrm{Mob}}
 =([0,2\pi]\times(-1,1))/
            ((0,t)\sim(2\pi,-t)).
\]
Wzdłuż środkowego okręgu $t=0$
wektor normalny $\partial_t$ zmienia
znak po obejściu okręgu; jego wiązka
normalna ma dokładnie takie samo
sklejenie $(0,v)\sim(2\pi,-v)$.
Nie istnieje globalne niezerowe
pole normalne: jego współczynnik
$a(\theta)$ musiałby być ciągły,
niezerowy i spełniać
$a(2\pi)=-a(0)$, co przeczy
twierdzeniu o wartości pośredniej.
\end{przyklad}

\begin{przyklad}[Przekątna i styczne do $M$]
\label{prz:normalna-diagonalna}
Dla $\Delta=\{(x,x):x\in M\}\subset M\times M$
mamy $T_{(x,x)}(M\times M)=T_xM\oplus T_xM$
i $T_{(x,x)}\Delta=\{(u,u):u\in T_xM\}$.
Odwzorowanie
\[
 [(u,v)]\longmapsto v-u
\]
jest izomorfizmem $\nu(\Delta,M\times M)\cong TM$:
jest surjektywne, jego jądrem przed
przejściem do ilorazu jest właśnie
przekątna w $T_xM\oplus T_xM$,
a wzór jest gładki w mapach.
Wiązka normalna może więc mieć
taką samą topologię jak wiązka styczna
i wcale nie musi być trywialna.
\end{przyklad}

\section{Otoczenie tubularne}
\label{sec:otoczenie-tubularne}

Zapiszmy
$D_\varepsilon\nu
=\{(x,v)\in\nu(S,M):|v|_g<\varepsilon(x)\}$
dla dodatniej funkcji $\varepsilon:S\to(0,\infty)$.
Przez Stwierdzenie~\ref{prop:normalna-rozszczepienie}
traktujemy tu $v$ jako wektor prostopadły
w $T_xM$.

\begin{twierdzenie}[Otoczenie tubularne]
\label{thm:otoczenie-tubularne}
Ustalmy dowolną metrykę riemannowską $g$ na $M$.
Jeśli $S\subset M$ jest domkniętą podrozmaitością
zanurzoną bez brzegu, istnieje dodatnia gładka
funkcja $\varepsilon$ i dyfeomorfizm
\[
 E:D_\varepsilon\nu(S,M)\longrightarrow U\subset M,
 \qquad E(x,v)=\exp_x^M(v),
\]
gdzie $U$ jest otwartym otoczeniem $S$,
a $E(x,0)=x$. Jeśli $S$ jest zwarta,
można wybrać stałe $\varepsilon>0$.
Każda zanurzona podrozmaitość, także
niedomknięta, ma przynajmniej pewne
otoczenie dyfeomorficzne z otoczeniem
przekroju zerowego swojej wiązki normalnej.
\end{twierdzenie}
\begin{proof}
\emph{Krok lokalny.} Przez cały dowód zachowujemy
ustaloną metrykę $g$; nie zakładamy jej zupełności.
Odwzorowanie
$E(x,v)=\exp_x(v)$ jest określone i gładkie
w pewnym otoczeniu przekroju zerowego
w $(TS)^\perp$. W punkcie $(x,0)$
przestrzeń styczna wiązki rozkłada się
kanonicznie na $T_xS\oplus(T_xS)^\perp$.
Jeżeli zmieniamy $x$ w $S$ przy $v=0$,
$E(x,0)=x$, więc różniczka w pierwszym
kierunku to $u$. Jeśli zmieniamy tylko
$v$, własność
$\dd(\exp_x)_0=\id_{T_xM}$ daje $w$.
W konsekwencji
\begin{equation}\label{eq:tuba-rozniczka}
 \dd E_{(x,0)}(u,w)=u+w.
\end{equation}
Jest to izomorfizm
$T_xS\oplus(T_xS)^\perp\to T_xM$.
Twierdzenie o odwzorowaniu odwrotnym
daje dyfeomorfizm na otoczeniu
każdego $(x,0)$.

\emph{Jednolita szerokość nad zwartym
fragmentem.} Niech $A\subset S$ będzie
zwarty. Istnieje $\delta_A>0$ takie, że
$E$ jest określone, jest lokalnym
dyfeomorfizmem i jest różnowartościowe
na wektorach długości mniejszej niż
$\delta_A$, których punkty bazowe
leżą w $A$. Istotne jest uzasadnienie
różnowartościowości. Gdyby dla dowolnie
małych długości zachodziła kolizja
$E(x_j,v_j)=E(y_j,w_j)$ z $x_j,y_j\in A$
i $(x_j,v_j)\ne(y_j,w_j)$, to po
przejściu do podciągu
$x_j\to x$, $y_j\to y$, $v_j,w_j\to0$.
Granica równości daje $x=y$.
Obie pary leżałyby wtedy ostatecznie
w jednym otoczeniu $(x,0)$, na którym
$E$ jest różnowartościowe na mocy
twierdzenia odwrotnego: sprzeczność.
Gładkość $E$ i zwartość $A$ dają
jednocześnie jednolitą dziedzinę
określoności i lokalną odwracalność.
Gdy $S$ jest zwarta, bierzemy $A=S$.

\emph{Niezwarta podrozmaitość domknięta.}
Wybierzmy gładkie właściwe wyczerpanie
$\rho:M\to[0,\infty)$ z
Twierdzenia~\ref{thm:wyczerpanie-15}; jego dowód
nie korzysta z otoczeń tubularnych.
Zbiory $A_j=S\cap\{\rho\leq j\}$
są zwarte, gdyż $S$ jest domknięty.
Weźmy malejące liczby $\delta_j>0$
takie, że $E$ jest różnowartościowe
i lokalnie odwracalne nad $A_j$
przy $|v|<\delta_j$, a ponadto
\[
 |\rho(E(x,v))-\rho(x)|<\tfrac14
 \qquad(x\in A_j,\ |v|<\delta_j).
\]
Ostatni warunek wynika z ciągłości $\rho\circ E$
i zwartości $A_j$: wzdłuż przekroju zerowego
różnica jest równa zeru. Możemy zmniejszać kolejne
$\delta_j$, zachowując wszystkie te własności.
Dobierzmy gładką
dodatnią $\varepsilon$ tak, aby
\[
  \varepsilon(x)<\tfrac14
  \quad\text{oraz}\quad
  \varepsilon(x)<\delta_{j+3}
  \ \text{dla}\ x\in S,\
       j-1\leq\rho(x)\leq j.
\]
Istnienie takiej funkcji wynika
z lokalnie skończonego podziału
jedności: na każdym otwartym kawałku
nieco szerszym od jednego pierścienia
wybieramy dodatnią stałą mniejszą
od skończenie wielu sąsiednich
$\delta_j$ i składamy ją z wagami
podziału. Lokalna skończoność
gwarantuje gładkość i dodatniość.

Załóżmy $E(x,v)=E(y,w)$ dla par
z $D_\varepsilon\nu$ i wybierzmy $j$
z $j-1\leq\rho(x)\leq j$.
Oznaczmy wspólny obraz przez $z$. Każdy wektor
spełnia ograniczenie nad swoim pierścieniem, zatem
\[
 |\rho(x)-\rho(y)|
 \leq |\rho(x)-\rho(z)|+|\rho(z)-\rho(y)|<\tfrac12.
\]
W szczególności $x,y\in A_{j+2}$.
Punkt $y$ leży najwyżej w sąsiednim
pierścieniu o numerze $j+1$; dolna
granica $\rho(y)$ może należeć
do pierścienia $j-1$. Z monotoniczności
$\delta_\ell$ i podanych ograniczeń
dla $\varepsilon$ wynika w każdym
przypadku $|v|,|w|<\delta_{j+2}$.
Jednolita różnowartościowość nad
$A_{j+2}$ wymusza $(x,v)=(y,w)$.
Tak $E$ jest różnowartościowym
lokalnym dyfeomorfizmem, a zatem
dyfeomorfizmem na otwarty obraz $U$.
Wykorzystanie $\rho$ działa
jednocześnie na wszystkich
składowych $M$.

Jeżeli $S$ nie jest domknięta w $M$,
wybierzmy dla każdego $x\in S$
mapę plasterkową o dziedzinie $V_x$.
Zbiór $S\cap V_x$ jest domknięty
w $V_x$, bo w tej mapie jest
przecięciem z domkniętą
podprzestrzenią $\R^k\times\{0\}$.
Niech $V=\bigcup_{x\in S}V_x$.
Domkniętość jest własnością
lokalną: dla każdego $y\in V$
wybieramy $V_x$ zawierający $y$,
a w nim $S\cap V_x$ jest domknięty.
Stąd $S$ jest domknięta w otwartym
otoczeniu $V\subset M$.
Stosujemy udowodnioną część
twierdzenia do zanurzenia $S\subset V$.
Ponieważ $TV|_S=TM|_S$, obie
wiązki normalne są tą samą wiązką.
Obraz uzyskanego dyfeomorfizmu
jest również otwarty w $M$. Używając metryki $g|_V$,
zachowujemy nawet opis przez $\exp^M$: na dostatecznie
krótkich geodezyjnych pozostających w $V$ oba odwzorowania
wykładnicze są identyczne.
\end{proof}

\begin{figure}[htbp]
\centering
\begin{tikzpicture}[>=Latex,font=\small,line cap=round,scale=1.03,
 declare function={tubeheight(\t)=.35*sin(deg(1.1*\t))-.65;
 tubeslope(\t)=.385*cos(deg(1.1*\t));
 tubenx(\t)=-tubeslope(\t)/sqrt(1+tubeslope(\t)^2);
 tubeny(\t)=1/sqrt(1+tubeslope(\t)^2);}]
 \fill[manifoldblue!9]
  plot[domain=-3.4:3.4,samples=100,variable=\t]
   ({\t+.42*tubenx(\t)},{tubeheight(\t)+.42*tubeny(\t)})
  -- plot[domain=3.4:-3.4,samples=100,variable=\t]
   ({\t-.42*tubenx(\t)},{tubeheight(\t)-.42*tubeny(\t)}) -- cycle;
 \foreach \r in {-.42,.42}
  \draw[manifoldblue!55,thin]
   plot[domain=-3.4:3.4,samples=100,variable=\t]
    ({\t+\r*tubenx(\t)},{tubeheight(\t)+\r*tubeny(\t)});
 \draw[manifoldblue,very thick]
  plot[domain=-3.4:3.4,samples=100,variable=\t] (\t,{tubeheight(\t)});
 \foreach \t in {-2.5,-1.1,.7,2.3}
  \draw[->,manifoldgold!85!black,thick]
   (\t,{tubeheight(\t)}) -- ({\t+.34*tubenx(\t)},{tubeheight(\t)+.34*tubeny(\t)});
 \fill[manifoldred] (.7,{tubeheight(.7)}) circle (2pt)
   node[below right,manifoldred] {$x$};
 \node[above,manifoldgold!65!black]
   at ({.7+.34*tubenx(.7)},{tubeheight(.7)+.34*tubeny(.7)}) {$E(x,v)$};
 \node[below left,manifoldblue] at (-3.4,{tubeheight(-3.4)}) {$S$};
 \node[anchor=west,manifoldblue!70!black] at (1.2,-1.4)
   {$U\simeq D_\varepsilon\nu$};
\end{tikzpicture}
\caption{Otoczenie tubularne krzywej $S$ (ciemnoniebieska).
W metryce euklidesowej złote odcinki są prostopadłe do
stycznych krzywej, a $E(x,v)=x+v$ umieszcza ich końce
w jasnym pasie $U$. Brzegi pasa są odsunięte wzdłuż normalnych.}
\label{fig:tubularna-15}
\end{figure}

\begin{wniosek}[Cofnięcie otoczenia na podrozmaitość]
\label{cor:tuba-retrakcja}
W otoczeniu tubularnym istnieje gładkie
rzutowanie na podstawę
$r= \pr_S\circ E^{-1}:U\to S$.
Zanurzenie $S\hookrightarrow U$ jest
równoważnością homotopijną.
W szczególności $H_*(U;G)\cong H_*(S;G)$
dla dowolnej abelowej grupy współczynników $G$.
\end{wniosek}
\begin{proof}
Na dysku wiązki definiujemy
$h_t(x,v)=(x,(1-t)v)$, $0\leq t\leq1$.
Ponieważ $|(1-t)v|\leq|v|<\varepsilon(x)$,
jest to dobrze określona homotopia
od identyczności do rzutu na przekrój
zerowy, stale równa identyczności
na tym przekroju. Przeniesienie przez
$E$ daje silną retrakcję deformacyjną
$U$ na $S$. Niezmienniczość homotopijna
homologii była udowodniona w
Rozdziale~\ref{ch:homologia-przestrzeni}.
\end{proof}

\begin{przyklad}[Okrąg i sfera]
W $\R^2$ dla okręgu jednostkowego
$E(e^{i\theta},t)=(1+t)e^{i\theta}$
przy $|t|<1$. Obrazem jest pierścień
$0<|z|<2$, a $r(z)=z/|z|$.
W $\R^3$ analogicznie
$E(x,t)=(1+t)x$ dla $x\in S^2$;
mały dysk wiązki jest powłoką
sferyczną. Na wstędze Möbiusa
otoczenie środkowego okręgu
jest skręconą wiązką dyskową,
nie iloczynem $S^1\times(-\varepsilon,\varepsilon)$.
\end{przyklad}

\section{Obramowanie i dane potrzebne do chirurgii}
\label{sec:obramowanie-normalne}

\begin{definicja}[Obramowanie]
\index{obramowanie@Obramowanie}
\label{def:obramowanie-normalne}
\emph{Obramowaniem normalnym} (ang.\ \emph{framing})
zanurzenia $S^k\subset M^n$ nazywamy
uporządkowaną bazę gładkich przekrojów
$(n_1,\ldots,n_{n-k})$ wiązki $\nu(S,M)$
w każdym włóknie; równoważnie,
izomorfizm wiązek
$\nu(S,M)\cong S\times\R^{n-k}$.
\end{definicja}

Istnienie obramowania to silniejszy
warunek niż orientowalność wiązki:
obramowanie daje orientację, lecz
orientacja nie zawsze daje globalną
bazę. Po wyborze obramowania
twierdzenie tubularne przy zwartej
$S$ daje zanurzenie
\begin{equation}\label{eq:obramowany-dysk}
 S\times D^{\,n-k}_{\varepsilon}
 \longrightarrow M,\qquad
 (x,a_1,\ldots,a_{n-k})
 \longmapsto
 \exp_x\!\left(\sum_j a_jn_j(x)\right).
\end{equation}
Tutaj $D^r_\varepsilon=\{a\in\R^r:|a|\leq\varepsilon\}$
jest \emph{domkniętym} dyskiem, w odróżnieniu od otwartego
$D_\varepsilon\nu$ z twierdzenia tubularnego.
Istotnie, na zwartej $S$ istnieje $C>0$ takie, że
$|\sum a_jn_j(x)|_g\leq C|a|$. Wybieramy
$C\varepsilon$ ściśle mniejsze od jednolitego promienia
tubularnego. Wówczas wzór jest zanurzeniem nawet
na otoczeniu tego domkniętego dysku. Obramowanie
nie musi być ortonormalne.
Bez obramowania dziedziną jest
dysk \emph{wiązki} normalnej,
a nie globalny iloczyn.

\begin{przyklad}[Równik w sferze]
Włóżmy $S^k\subset S^n\subset\R^{n+1}$
jako punkty $(x,0)$ w rozkładzie
$\R^{n+1}=\R^{k+1}\oplus\R^{n-k}$.
Stałe wektory $e_{k+2},\ldots,e_{n+1}$
są w każdym punkcie $(x,0)$ styczne
do $S^n$, lecz prostopadłe do $T_xS^k$.
Tworzą więc obramowanie normalne.
W kontraście środkowy okrąg wstęgi
Möbiusa z Przykładu~\ref{prz:normalna-mobius}
nie dopuszcza nawet jednego
niezerowego normalnego przekroju.
\end{przyklad}

W teorii chirurgii wybiera się
często zanurzenie sfery wraz z
obramowaniem. Wzór~\eqref{eq:obramowany-dysk}
mówi wtedy dokładnie, \emph{które}
otoczenie $S^k\times D^{n-k}$
ma zostać usunięte. Różne obramowania
mogą dawać różne sklejenia po usunięciu;
samo wskazanie podrozmaitości
nie zawiera całej potrzebnej informacji.

\section{Kołnierz brzegu i podwojenie}
\label{sec:kolnierz-15}

\begin{definicja}[Kołnierz]
\index{kolnierz@Kołnierz}
Jeżeli $W$ jest rozmaitością z brzegiem,
\emph{kołnierzem brzegu} nazywamy zanurzenie
\[
 c:\partial W\times[0,1)\longrightarrow W,
 \qquad c(p,0)=p,
\]
którego obraz jest otwartym otoczeniem
$\partial W$ w $W$. Parametr $t$
oznacza ruch \emph{do wnętrza}.
\end{definicja}

\begin{twierdzenie}[Istnienie kołnierza]
\label{thm:kolnierz-15}
Każda gładka rozmaitość z brzegiem
ma kołnierz. Jeśli $\partial W$ jest
zwarty, kołnierz można przed
przeskalowaniem zapisać na
$\partial W\times[0,\varepsilon)$
z jedną stałą $\varepsilon>0$.
\end{twierdzenie}
\begin{proof}
W mapie półprzestrzennej
$(u^1,\ldots,u^{n-1},r)$, $r\geq0$,
wektor $\partial_r$ wskazuje do
wnętrza w punktach brzegu. Wybierzmy
lokalnie skończone pokrycie brzegu
takimi mapami i podział jedności
na niewielkim otoczeniu brzegu.
Ważona suma odpowiednich pól
$\partial_r$ jest gładkim polem
$V$, którego składowa skierowana
do wnętrza jest dodatnia:
każda dodatnia kombinacja
wektorów skierowanych do wnętrza
nadal jest skierowana do wnętrza.

Lokalny przepływ $\Phi_t$ tego pola przy brzegu
otrzymujemy, przedłużając jego współczynniki
w mapach półprzestrzennych przez $r=0$ i stosując
twierdzenie o równaniach różniczkowych. Jednoznaczność
spaja te lokalne rozwiązania, dopóki leżą w $W$.
Dla każdego $p\in\partial W$ istnieje otoczenie
$B_p$ w brzegu i czas $\tau_p>0$, dla których
$\Phi_t(q)$ jest określone w dziedzinie $V$
i leży we wnętrzu $W$ dla $q\in B_p$ i $0<t<\tau_p$.
Wynika to z dodatniej składowej pola w kierunku $r$
oraz ciągłej zależności rozwiązania od danych.

Za pomocą lokalnie skończonego podziału jedności
na brzegu wybieramy dodatnią gładką funkcję
$\varepsilon(p)$ tak małą, by te własności zachodziły
dla $0<t<\varepsilon(p)$. Dokładniej, dla podziału
$\{\chi_i\}$ podporządkowanego otoczeniom $B_i$
można wziąć $\varepsilon=\sum_i\chi_i\tau_i/2$:
przy ustalonym $p$ jest to średnia skończenie wielu
dopuszczalnych czasów, więc jest mniejsza od
największego z nich, także dopuszczalnego przy $p$.

Odwzorowanie $C(p,t)=\Phi_t(p)$ ma różniczkę
\[
 \dd C_{(p,t)}(u,a)=\dd\Phi_t{}_p(u+aV_p).
\]
Jest ona izomorfizmem: pole $V_p$ jest poprzeczne
do $T_p\partial W$, a różniczka lokalnego przepływu
jest odwracalna. Przy $t=0$ stosujemy twierdzenie
odwrotne w mapach półprzestrzennych; skierowanie
$V$ do wnętrza zapewnia właściwą stronę brzegu.
Tak $C$ jest lokalnym dyfeomorfizmem na dziedzinie
$\{(p,t):0\leq t<\varepsilon(p)\}$.

Globalna różnowartościowość wynika bezpośrednio
z jednoznaczności trajektorii. Jeżeli
$\Phi_t(p)=\Phi_s(q)$ i $t\geq s$, to cofając
obie trajektorie o czas $s$, dostajemy
$\Phi_{t-s}(p)=q$. Dla $t>s$ lewa strona
leży we wnętrzu, a $q$ na brzegu, co jest
niemożliwe. Zatem $t=s$ i $p=q$.
Różnowartościowy lokalny dyfeomorfizm $C$ jest
dyfeomorfizmem na otwarte otoczenie brzegu.
Wreszcie
\[
 c(p,s)=C(p,s\varepsilon(p)),\qquad 0\leq s<1,
\]
daje kołnierz z definicji. Przy zwartym brzegu
minimum dodatniej funkcji $\varepsilon$ jest dodatnie,
więc przed przeskalowaniem można użyć stałej szerokości.
\end{proof}

\begin{figure}[htbp]
\centering
\begin{tikzpicture}[>=Latex,font=\small,line cap=round]
 \fill[manifoldblue!8] (-3,0) rectangle (3,1.35);
 \fill[manifoldgreen!8] (-3,1.35) rectangle (3,2.15);
 \draw[manifoldblue,very thick] (-3,0)--(3,0);
 \draw[manifoldblue!45] (-3,1.35)--(3,1.35);
 \foreach \x in {-2.4,-1.4,-.4,.6,1.6,2.6}
   \draw[->,manifoldgold!85!black,thick]
      (\x,0.12)--(\x,1.02);
 \fill[manifoldred] (-1.4,0) circle (2pt);
 \node[anchor=north,manifoldred] at (-1.4,-.08) {$p\in\partial W$};
 \node[anchor=west,manifoldblue] at (1.6,-.32) {$\partial W$};
 \draw[->,manifoldred,thick] (-1.4,.12)--(-1.4,1.02);
 \draw[manifoldred,thin] (-1.4,1.1)--(-.8,1.7);
 \node[anchor=west,manifoldred] at (-.7,1.75)
       {$t\mapsto c(p,t)$};
 \node[anchor=west,manifoldgreen!60!black] at (-2.8,1.75)
       {$\operatorname{int}W$};
\end{tikzpicture}
\caption{Kołnierz zamienia całe otoczenie brzegu
w produkt $\partial W\times[0,1)$.
Współrzędna $t$ rośnie do wnętrza rozmaitości.}
\label{fig:kolnierz-15}
\end{figure}

\begin{przyklad}[Dysk i walec]
Dla dysku jednostkowego $D^2$
kołnierz można zapisać jako
$c(p,t)=(1-t)p$ dla $p\in S^1$
i $0\leq t<\varepsilon<1$.
Dla $W=S^1\times[0,1]$ brzeg
ma dwie składowe: przy $0$
parametr idzie w górę, a przy
$1$ w dół. Jest to konkretny
powód, dla którego orientacje
dwóch składowych brzegu walca
są przeciwne.
\end{przyklad}

\begin{wniosek}[Podwojenie rozmaitości]
\label{cor:podwojenie-15}
Dwie kopie $W_+$ i $W_-$ rozmaitości
z brzegiem można skleić wzdłuż
identyczności $\partial W$ tak,
aby otrzymać gładką rozmaitość
bez brzegu
$DW=W_+\cup_{\partial W}W_-$.
Jeżeli $W$ jest zwarta, $DW$
też jest zwarta.
\end{wniosek}
\begin{proof}
Poza brzegiem korzystamy z
istniejących map gładkich obu kopii.
W pobliżu sklejonego brzegu
używamy kołnierzy i współrzędnej
$s\in(-1,1)$: dla $s\geq0$
punkt $(p,s)$ oznacza $c_+(p,s)$,
a dla $s\leq0$ oznacza
$c_-(p,-s)$. Wspólne przejścia
między takimi mapami mają postać
$(p,s)\mapsto(\psi\circ\varphi^{-1}(p),s)$,
więc są gładkie. Z kolei przejścia
do map pochodzących z wnętrza
odbywają się dla $s\ne0$ i także
są gładkie. Otrzymujemy atlas
bez półprzestrzennych punktów
brzegowych. Sprawdźmy też warunki topologiczne.
Różne punkty szwu rozdzielamy rozłącznymi otoczeniami
w $\partial W$ i ich kołnierzami. Punkt wnętrza
można oddzielić od szwu, zmniejszając otoczenie
tak, by jego domknięcie nie dochodziło do brzegu;
następnie dobieramy odpowiednio małe otoczenie
drugiego punktu. Punkty wnętrz rozdziela własność
Hausdorffa w kopiach $W$. Iloraz jest więc Hausdorffa.
Przeliczalne atlasy obu wnętrz i brzegu wraz
z przedziałami o wymiernych końcach dają przeliczalną
bazę ilorazu. Jeśli $W$ jest zwarta, ciągły obraz
zwartej sumy rozłącznej $W_+\sqcup W_-$ jest zwarty.
\end{proof}

\section{Izotopie i zgodność otoczeń}
\label{sec:izotopie-otoczen-15}

Izotopia to gładka rodzina zanurzeń
$i_t:S\hookrightarrow M$, $0\leq t\leq1$.
Rozdział~\ref{ch:pola} zawiera
już dowód twierdzenia o rozszerzeniu
izotopii~\ref{thm:rozszerzenie-izotopii}
dla zwartej $S$ bez brzegu:
istnieje izotopia otoczenia
$h_t:M\to M$ z $h_0=\id$
i $h_t\circ i_0=i_t$.
Nie powtarzamy tamtego dowodu;
wykorzystamy jego wnioski geometryczne.

\begin{wniosek}[Otoczenia i dopełnienia izotopijnych zanurzeń]
\label{cor:izotopia-dopelnienia}
Jeśli $S$ jest zwarta bez brzegu,
a $i_t:S\hookrightarrow M$ jest
gładką izotopią, to
\[
 M\setminus i_0(S)\cong M\setminus i_1(S)
 \quad\text{oraz}\quad
 \nu(i_0(S),M)\cong \nu(i_1(S),M)
\]
jako odpowiednio rozmaitości
i wiązki nad $S$. Można wybrać
otoczenia tubularne obu obrazów
odpowiadające sobie przez $h_1$.
\end{wniosek}
\begin{proof}
Twierdzenie~\ref{thm:rozszerzenie-izotopii}
daje dyfeomorfizm $h_1$ z
$h_1(i_0(S))=i_1(S)$.
Jego ograniczenie do dopełnień
ma odwrotność $h_1^{-1}$.
Różniczka $\dd h_1$ wysyła
$T(i_0(S))$ na $T(i_1(S))$,
a więc przechodzi do
izomorfizmu ilorazów
$TM|_{i_0(S)}/T(i_0(S))
\to TM|_{i_1(S)}/T(i_1(S))$.
Jeśli $E_0:D_\varepsilon\nu_0\to U_0$
jest otoczeniem tubularnym,
to złożenie $h_1\circ E_0$
opisuje otoczenie $h_1(U_0)$
przez tę samą wiązkę; po
uzyskanym izomorfizmie $\nu_0\cong\nu_1$
jest to opis przez $\nu_1$. Precyzyjniej, jeśli $A$
oznacza ten izomorfizm, nowym odwzorowaniem tubularnym jest
$E_1=h_1\circ E_0\circ A^{-1}$. Nie musi ono być
odwzorowaniem wykładniczym dla początkowej metryki na $M$.
Jest nim dla metryki przeniesionej przez $h_1$,
ponieważ izometria przenosi geodezyjne i normalne.
\end{proof}

\begin{uwaga}[Co właściwie jest zachowane?]
Wniosek nie mówi, że dowolne dwa
zanurzenia tej samej abstrakcyjnej
sfery mają dyfeomorficzne dopełnienia.
Wymagana jest \emph{izotopia zanurzeń}
w ustalonym $M$. W szczególności
różne węzły w $\R^3$ mogą mieć
różne dopełnienia; nie należy
utożsamiać deformacji abstrakcyjnego
okręgu z izotopią w przestrzeni
otaczającej.
\end{uwaga}

Kołnierze tej samej zwartej
składowej brzegu można także
porównywać: z kołnierza $c_j$
otrzymujemy pole
$(c_j)_*(\partial_t)$ w pobliżu
brzegu. Oba pola są skierowane
do wnętrza. Ich interpolacja
$(1-s)V_0+sV_1$ nadal jest
skierowana do wnętrza;
jej przepływ daje, po wyborze
wspólnego małego czasu na
zwartym brzegu, rodzinę
kołnierzy łączącą oba opisy.
Oznaczmy tę rodzinę przez $C_s(p,t)$. Aby otrzymać
dyfeomorfizm, a nie tylko rodzinę opisów, definiujemy
w jej obrazie zależne od $s$ pole
$Y_s(C_s(p,t))=\partial_s C_s(p,t)$.
Ponieważ $C_s(p,0)=p$, pole to znika na brzegu.
Wybieramy gładkie odcięcie równe $1$ na mniejszym
kołnierzu i równe $0$ przed zewnętrznym końcem większego.
Zwartość brzegu i przedziału parametrów pozwala dobrać
wspólną szerokość i zwarty nośnik dla całej rodziny.
Przepływ powstałego pola jest izotopią $H_s$ rozmaitości
$W$, ustaloną na brzegu; jednoznaczność rozwiązań daje
$H_s(C_0(p,t))=C_s(p,t)$ w mniejszym kołnierzu.
Stosując $H_1$ na obu kopiach $W$, otrzymujemy
dyfeomorfizm odpowiednich podwojeń: przy szwie
współrzędne $(p,\text{czas ze znakiem})$ są zachowane.
Wyjaśnia to niezależność gładkiego podwojenia, z dokładnością
do dyfeomorfizmu, od wyboru kołnierza przy zwartym brzegu.

\section{Funkcje wyczerpujące i regularne poziomice}
\label{sec:wyczerpujace-15}

\begin{definicja}[Funkcja wyczerpująca]
\index{funkcja wyczerpujaca@Funkcja wyczerpująca}
\label{def:wyczerpujaca-15}
Funkcja ciągła $\rho:M\to[0,\infty)$ jest
\emph{wyczerpująca}, jeśli jest właściwa,
czyli $\rho^{-1}([0,a])$ jest zwarty
dla każdego $a\geq0$.
Nie wymaga się, by $\rho$ była
funkcją Morse'a.
\end{definicja}

\begin{twierdzenie}[Istnienie gładkiego wyczerpania]
\label{thm:wyczerpanie-15}
Każda rozmaitość gładka z przeliczalną
bazą, również z brzegiem, ma gładką
nieujemną funkcję wyczerpującą.
Na rozmaitości bez brzegu istnieje
ciąg dowolnie dużych wartości regularnych
$a_j\to\infty$; zbiory
$\rho^{-1}((-\infty,a_j])$ są wtedy
zwartymi rozmaitościami z gładkim
brzegiem $\rho^{-1}(a_j)$.
\end{twierdzenie}
\begin{proof}
Z lokalnej zwartości i przeliczalnej
bazy budujemy zwarte zbiory
$K_1\subset\operatorname{int}K_2
\subset K_2\subset\operatorname{int}K_3
\subset\cdots$ o sumie równej $M$.
W tym celu wybieramy przeliczalne pokrycie
$\{B_j\}$ zbiorami otwartymi o zwartych domknięciach.
Mając $K_j$, pokrywamy zwarty zbiór
$K_j\cup\overline{B_{j+1}}$ skończenie wieloma
względnie zwartymi otwartymi otoczeniami i za
$K_{j+1}$ bierzemy sumę ich domknięć. Zaczynamy od
$K_1\supset\overline{B_1}$; konstrukcja zapewnia
zarówno zawieranie we wnętrzach, jak i wyczerpanie $M$.

Dla każdego $j$ skonstruujmy
gładką $\psi_j:M\to[0,1]$
równą $0$ w otoczeniu $K_j$
i równą $1$ poza
$\operatorname{int}K_{j+1}$.
Istnienie wynika z gładkiego
lematu o funkcji odcinającej:
zbiory $K_j$ i
$M\setminus\operatorname{int}K_{j+1}$
są rozłączne i domknięte, a
lokalnie skończony podział
jedności oddziela je gładko.
Połóżmy
\[
 \rho(x)=\sum_{j=1}^{\infty}\psi_j(x).
\]
Wybierzmy $m$ tak, by $x\in\operatorname{int}K_m$.
Wszystkie $\psi_j$ dla $j\geq m$ są równe zeru
na tym samym otoczeniu $\operatorname{int}K_m$;
suma jest lokalnie skończona,
zatem $\rho$ jest gładka.
Jeśli $x\notin K_{N+1}$,
to dla $j\leq N$ mamy
$x\notin\operatorname{int}K_{j+1}$,
więc $\psi_j(x)=1$ i $\rho(x)\geq N$.
Stąd $\{\rho\leq a\}$ zawiera się
w $K_{N+1}$ dla $N>a$.
Jest domknięty, zatem zwarty;
to dowodzi właściwości.

Dla rozmaitości bez brzegu
twierdzenie Sarda~\ref{thm:sard}
daje wartości regularne poza
zbiorem miary zero. W każdym
przedziale $(j,j+1)$ można
wybrać wartość regularną $a_j$.
W pobliżu każdego punktu jej
poziomicy twierdzenie o submersji
zapisuje $\rho$ jako jedną
ze współrzędnych, więc podpoziomica
ma lokalny model półprzestrzeni.
Jej brzeg jest właśnie
$\rho^{-1}(a_j)$.
\end{proof}

\begin{twierdzenie}[Właściwa funkcja Morse'a]
\label{thm:wlasciwa-morse-15}
Na każdej gładkiej rozmaitości bez brzegu
istnieje nieujemna właściwa funkcja Morse'a $f:M\to\R$.
Każda jej podpoziomica $\{f\leq a\}$ jest zwarta.
Jej punkty krytyczne mogą być nieskończenie liczne,
ale w każdej podpoziomicy jest ich skończenie wiele.
\end{twierdzenie}
\begin{proof}
Zaczynamy od gładkiej właściwej
funkcji $\rho$ z
Twierdzenia~\ref{thm:wyczerpanie-15}.
Wybierzmy zwarte
$K_1\subset\operatorname{int}K_2
\subset K_2\subset\cdots$
o sumie $M$; dla wygody
$K_0=\varnothing$. Połóżmy $f_0=\rho$.
Skonstruujemy funkcje
$f_j=f_{j-1}+h_j$ tak, aby $f_j$
była Morse'a na otoczeniu $K_j$,
a przyszłe poprawki nie niszczyły
tej własności.

Pokryjmy $K_j$ skończoną liczbą
mniejszych obszarów map. W większych
obszarach wybierzmy funkcje
odcinające $\beta_i$, równe $1$
na mniejszych. Funkcje
$\beta_i x_i^1,\ldots,\beta_i x_i^n$,
przedłużone przez zero, mają
różniczki rozpinające $T_x^*M$
dla każdego $x\in K_j$:
w jednym z mniejszych obszarów
są po prostu współrzędnymi.
Zatem dla rodziny
\[
 f_{j,a}=f_{j-1}
       +\sum_{i,\ell}a_{i\ell}\beta_i x_i^\ell
\]
parametryczna różniczka
$(x,a)\mapsto\dd f_{j,a}(x)$
jest transwersalna do przekroju
zerowego nad otoczeniem $K_j$.
Dowód transwersalności
parametrycznej z
Twierdzenia~\ref{thm:transwersalnosc-param}
pozwala wybrać \emph{dowolnie
małe} parametry $a$, dla których
$f_{j,a}$ jest Morse'a na
otoczeniu $K_j$. Niech $h_j$
będzie otrzymaną poprawką.
Jest sumą skończenie wielu
funkcji o zwartych nośnikach.

Uściślijmy wybór wielkości poprawek. Ustalmy
metrykę riemannowską i jej połączenie $\nabla$ oraz normy
\[
 \|h\|_{C^r(K)}=\max_{0\leq\ell\leq r}
       \sup_{x\in K}|\nabla^\ell h(x)|.
\]
Na zbiorach zwartych są one równoważne normom
w skończonych układach map; rosną wraz z $K$ i $r$.
Normy na pustym $K_0$ przyjmujemy równe zeru.
Skoro $f_i$ jest Morse'a na
otoczeniu zwartego $K_i$,
istnieje $\eta_i>0$ takie,
że perturbacja o normie
$C^2(K_i)$ mniejszej od
$\eta_i$ zachowuje
niezdegenerowanie wszystkich
punktów krytycznych na $K_i$.
Wynika to z otwartości
transwersalności $\dd f_i$
do przekroju zerowego
na zwartym zbiorze.
W kroku $j$ wybieramy małe
parametry tak, aby dodatkowo
\[
 \|h_j\|_\infty<2^{-j},\qquad
 \|h_j\|_{C^j(K_{j-1})}<2^{-j},\qquad
 \|h_j\|_{C^2(K_i)}
          <2^{-j}\eta_i\quad(i<j).
\]
Każdy krok narzuca tylko
skończenie wiele warunków
otwartych przy $a=0$,
więc można je pogodzić
z wyborem dobrego $a$.

Szereg $f=\rho+\sum_{j\geq1}h_j$
zbiega gładko na każdym zwartym
zbiorze: dla ustalonego $K_i$
i rzędu pochodnej $r$ wszystkie
wyrazy o $j>\max(i,r)$ mają normę
$C^r(K_i)$ ograniczoną przez
$2^{-j}$. Ponadto
$|f-\rho|\leq\sum_j2^{-j}=1$.
Zatem przeciwobraz zwartego
przedziału $[-A,A]$ przez $f$
leży w zwartej podpoziomicy
$\{\rho\leq A+1\}$ i jest
domknięty, czyli zwarty.
Funkcja $f$ jest właściwa.
Na $K_i$ różnica $f-f_i$
ma normę $C^2$ mniejszą
od $\eta_i\sum_{j>i}2^{-j}
<\eta_i$; funkcja $f$
pozostaje tam Morse'a.
Ponieważ $K_i$ wyczerpują
$M$, jest Morse'a wszędzie. Ponieważ $f\geq-1$,
zastępujemy ją na końcu przez $f+1$. Otrzymana funkcja
jest nieujemna, nadal właściwa i Morse'a, a jej
podpoziomice są przeciwobrazami zwartych przedziałów.
Jej punkty krytyczne są izolowane;
ich zbiór w zwartej podpoziomicy
jest domknięty i dlatego skończony.
\end{proof}

\begin{stwierdzenie}[Metryka zupełna]
\label{prop:metryka-zupelna-wyczerpanie}
Każda rozmaitość gładka bez brzegu
ma zupełną metrykę riemannowską.
Ten sam wniosek w sensie zupełnej
odległości metrycznej zachodzi
dla rozmaitości z brzegiem.
\end{stwierdzenie}
\begin{proof}
Wybierzmy dowolną metrykę $g$
i gładkie właściwe $\rho$
z Twierdzenia~\ref{thm:wyczerpanie-15}.
Niech
$\lambda=\sqrt{1+|\nabla_g\rho|_g^2}$
i $g'=\lambda^2g$.
Ponieważ
$|\nabla_{g'}\rho|_{g'}
=|\nabla_g\rho|_g/\lambda\leq1$,
całkowanie po dowolnej kawałkami
gładkiej krzywej $\gamma$ daje
\[
 |\rho(\gamma(1))-\rho(\gamma(0))|
 \leq L_{g'}(\gamma).
\]
Po wzięciu kresu dolnego długości
$|\rho(x)-\rho(y)|\leq d_{g'}(x,y)$.
Niech $(x_j)$ będzie ciągiem Cauchy'ego
w odległości $d_{g'}$. Jego odległości
od ustalonego $x_1$ są ograniczone,
więc wartości $\rho(x_j)$ też są
ograniczone. Cały ciąg leży
w pewnej zwartej podpoziomicy
$\{\rho\leq A\}$, zatem ma podciąg
zbieżny. Ponieważ $(x_j)$ jest
Cauchy'ego i topologia $d_{g'}$
zgadza się z topologią rozmaitości,
cały ciąg zbiega. Ten sam argument
dla krzywych leżących w rozmaitości
z brzegiem daje jej zupełność
jako przestrzeni metrycznej. Zgodność topologii
przy brzegu wynika z lokalnego porównania metryki
z euklidesową w mniejszych wypukłych półkulach map.
Przy niespójnym $M$ zupełność riemannowską rozumiemy
składowa po składowej, ponieważ między różnymi składowymi
nie ma krzywych. Jeśli potrzebna jest jedna skończona
odległość na całym $M$, bierzemy $\min(d_{g'},1)$
na każdej składowej i odległość $1$ pomiędzy nimi.
Jest ona zupełna i wyznacza topologię $M$.
\end{proof}

\begin{przyklad}
Na $\R^n$ funkcja $\rho(x)=|x|^2$
jest wyczerpująca, bo
$\{\rho\leq a\}$ to domknięta
kula promienia $\sqrt a$.
Na otwartym przedziale $(0,1)$
działa na przykład
$\rho(x)=x^{-2}+(1-x)^{-2}$:
rośnie bez ograniczenia przy
obu końcach. Funkcja $x^2$
na $(0,1)$ nie jest właściwa:
punkty mogą zbliżać się do
brzegu $0$ przy ograniczonych
wartościach $x^2$.
\end{przyklad}

\begin{twierdzenie}[Pas bez punktów krytycznych jest produktem]
\label{thm:pas-bez-krytycznych-15}
Niech $M$ będzie rozmaitością bez brzegu,
$f:M\to\R$ gładką funkcją
i właściwą, a $a<b$. Załóżmy,
że $f$ nie ma punktów krytycznych
w $f^{-1}([a,b])$. Wtedy
\[
 f^{-1}([a,b])\cong f^{-1}(a)\times[a,b],
 \qquad f\ \text{odpowiada rzutowi na }[a,b].
\]
Jeżeli poziomica $f^{-1}(a)$ jest pusta,
obie strony są puste.
\end{twierdzenie}
\begin{proof}
Pas $K=f^{-1}([a,b])$ jest zwarty,
bo $f$ jest właściwa. Dla $K=\varnothing$ teza
jest natychmiastowa. W przeciwnym razie na nim
$\nabla f\ne0$, więc pole
\[
 X=\frac{\nabla f}{|\nabla f|^2}
\]
jest określone w otoczeniu $K$.
Wybierzmy gładką funkcję
odcinającą $\chi$ o zwartym
nośniku w tym otoczeniu,
równą $1$ na $K$, i zastąpmy
$X$ przez globalne, zwartego
nośnika pole $\chi X$.
Jego przepływ $\Phi_t$
istnieje dla wszystkich
$t\in\R$ dzięki zwartości
nośnika. Dopóki trajektoria
pozostaje w $K$, mamy
\[
 \frac{\dd}{\dd t}f(\Phi_t(x))
 =\dd f_{\Phi_t(x)}(X)
 =g(\nabla f,\nabla f)/|\nabla f|^2=1.
\]
Startując z $x\in f^{-1}(a)$,
trajektoria nie może opuścić
$K$ przed czasem $b-a$:
wcześniejsze opuszczenie
wymagałoby wartości $f$
poza $[a,b]$, podczas gdy
na każdym wcześniejszym
odcinku $f(\Phi_t(x))=a+t$.
Zatem
$\Psi(x,s)=\Phi_{s-a}(x)$
dla $s\in[a,b]$ jest
dobrze określona i ma
$f(\Psi(x,s))=s$.
Analogicznie trajektoria startująca z $y\in K$
pozostaje w $K$ dla $a-f(y)\leq t\leq0$ i spełnia
$f(\Phi_t(y))=f(y)+t$. Dochodzi więc do poziomicy $a$;
jeśli ta jest pusta, nie mogło być żadnego $y\in K$.
Własność grupowa przepływu pokazuje, że odwrotnością
$\Psi$ jest
\[
 y\longmapsto
       \bigl(\Phi_{a-f(y)}(y),\,f(y)\bigr).
\]
Przepływ i $f$ są gładkie,
więc oba odwzorowania są gładkie,
także w mapach półprzestrzennych
przy $a$ i $b$. To daje
żądaną identyfikację.
\end{proof}

\begin{figure}[htbp]
\centering
\begin{tikzpicture}[>=Latex,font=\small,line cap=round]
 \fill[manifoldblue!7] (-2.4,.55) rectangle (2.4,2.25);
 \draw[manifoldblue!65,thick] (-2.4,.55)--(2.4,.55);
 \draw[manifoldblue!65,thick] (-2.4,2.25)--(2.4,2.25);
 \draw[manifoldblue!35,dashed] (-2.4,1.45)--(2.4,1.45);
 \foreach \x in {-1.5,-.4,.7,1.8}
  \draw[->,manifoldgold!80!black,thick] (\x,.75)--(\x,2.02);
 \node[anchor=west] at (2.5,.55) {$f=a$};
 \node[anchor=west] at (2.5,2.25) {$f=b$};
 \node[above,manifoldgold!70!black] at (0,2.5)
       {$X=\nabla f/|\nabla f|^2=\partial_s$};
\end{tikzpicture}
\caption{Przepływ pola $X$ przenosi punkt poziomicy
$f=a$ na poziomicę $f=a+t$. Gdy w zwartym pasie
nie ma punktów krytycznych, wszystkie poziomice
składają się w produkt. Rysunek przedstawia współrzędne
produktowe $(x,s)$ wyznaczone przez przepływ; w nich $f=s$.}
\label{fig:pas-bez-krytycznych-15}
\end{figure}

\begin{uwaga}[Związek z późniejszą teorią Morse'a]
Jeśli $\rho$ jest właściwa,
zmiany topologii jej podpoziomic
mogą zachodzić tylko przy
przechodzeniu przez wartości
krytyczne: pomiędzy nimi
Twierdzenie~\ref{thm:pas-bez-krytycznych-15}
identyfikuje pas z produktem.
Teoria Morse'a opisze
konkretną zmianę przy
niezdegenerowanym punkcie
krytycznym. Tu potrzebowaliśmy
jedynie pojęcia punktu
krytycznego z Rozdziału~\ref{ch:sard-transwersalnosc}.
\end{uwaga}

\par\smallskip
{\footnotesize\noindent\emph{Dalsza lektura:}
\href{https://ocw.mit.edu/courses/18-965-geometry-of-manifolds-fall-2004/resources/lecture21_22/}{wykłady MIT o mocnym twierdzeniu Whitneya},
\href{https://arxiv.org/abs/1907.10297}{omówienie dysku Whitneya w wykładach Benedettiego},
\href{https://ocw.mit.edu/courses/18-966-geometry-of-manifolds-spring-2007/resources/lect05/}{wykład MIT o otoczeniach tubularnych}.\par}
